% supp_sim.tex -- Supplementary Section S5: simulations of the scenes and of interventions; details.
% Paths in "% src:" comments are relative to the repository root.
% Table bodies between "% BEGIN AUTO" and "% END AUTO" are written by code/s6_scenes_tables.py; derived
% numbers: data/s6_scenes_numbers.json, data/s7_interventions_numbers.json and the files named.
\section{Simulations of the scenes and of interventions: details}\label{appS_sim:sec}

This section gives the details of Section~\ref{M-s6_scenes:sec} of the main text. Unless stated otherwise
$\beta=0.5\perday$, $\gamma=0.14\perday$, the contact kernel is the 1\,m kernel, there is one uniformly placed index
case, the room is occupied 24 hours a day, a `major' outbreak has an attack rate of at least 10\,\%, and $\Rhat$ is
the counted number of secondary cases of an index case that alone transmits. $\Ronebar$ and $\srad(\ngmone)$ are not
threshold quantities (Remark~\ref{M-s5_pair:rem:threshold}).

\subsection{The four scenes}\label{appS_sim:sec:scenes}

An epidemic in the office scene at $\Dz=1$ consists of about 5{,}900 events. Peak prevalence and peak day are the
maximum of $I/\Npop$ on a grid of $0.5$~day and its time, averaged over major outbreaks. The mean-field equations
are started from one infective spread uniformly over the cells, and from one infective placed in a single cell
(averaged over every third cell as seed); the extinction probabilities of \eqref{M-s6_scenes:eq:branching} are
computed by iteration from $s=0$. Fig.~\ref{s6_scenes:fig:curves} shows the epidemic curves of the four scenes
and Fig.~\ref{s6_scenes:fig:attack} the distribution of the attack rate over the same runs.
% src: data/bsc_sim/03_scenes.json: "office|D0=1|range1m".events_per_rep = 5919.5; code/bsc_sim/experiments/03_scenes.py (n_rep = 2000, T_MAX = 400 d); 03_scenes.json: n_unfinished = 0 in all 18 rows; code/s6_scenes_ode_pointseed.py

\begin{figure}[tbp]
\centering
\includegraphics[width=\textwidth]{fig5_3_scene_epidemic_curves_en}
\caption{Prevalence $I/\Npop$ against time in the four scenes (columns; $\Npop=100$, $150$, $80$, $310$) at
$\Dz=1$ (top) and $\Dz=10\cellday$ (bottom); 1\,m kernel, 2{,}000 exact simulations per panel, one uniformly placed
index case. Solid line: mean over major outbreaks (attack rate at least $10\,\%$) at each time, with the band between
the 5th and 95th percentiles; the dot marks its maximum. Dotted: mean over all runs. Dashed: mean-field equations
\eqref{M-s2_model:eq:meanfield} started from one infective spread uniformly over the cells. The values of
$P_{\mathrm{maj}}$ in the panels are those of Table~\ref{M-s6_scenes:tab:scenes}. The maximum of the mean curve (office, $\Dz=1$: $7.0\,\%$
on day $19.5$) is lower than the mean of the individual peaks in Table~\ref{M-s6_scenes:tab:scenes} ($12.7\,\%$)
because outbreaks peak at different times.}
% src: figures/fig5_3_scene_epidemic_curves_en.pdf (code/bsc_sim/experiments/21_figures_more.py <- data/bsc_sim/03_curves_*.npz);
%      data/s6_scenes_numbers.json: scenes."office|D0=1".mean_curve_major_peak = 0.0696, _day = 19.5
\label{s6_scenes:fig:curves}
\end{figure}

\begin{figure}[tbp]
\centering
\includegraphics[width=\textwidth]{fig5_4_attack_rate_distributions_en}
\caption{Distribution of the attack rate (fraction of the $\Npop$ people ever infected) over the 2{,}000 runs of
each panel of Fig.~\ref{s6_scenes:fig:curves}; bins of width $0.04$, logarithmic vertical axis. The dashed line is
the $10\,\%$ threshold that defines a major outbreak. At $\Dz=10\cellday$ the distributions are bimodal; at $\Dz=1$
they are not, except in the metro carriage.}
% src: figures/fig5_4_attack_rate_distributions_en.pdf (code/bsc_sim/experiments/21_figures_more.py <- data/bsc_sim/03_curves_*.npz)
\label{s6_scenes:fig:attack}
\end{figure}

The values of $\Rhat$ in Table~\ref{M-s6_scenes:tab:scenes} come from separate batches of 2{,}000 runs and are the
least precise estimates of this quantity in the article; for the office, samples of 100{,}000 to 400{,}000
runs give $2.265\pm0.004$, $3.793\pm0.006$ and $4.340\pm0.015$ at $\Dz=1$, $10$ and $100$
(Section~\ref{M-s5_pair:sec:mobility}), in agreement with $\Ronebar$.
% src: data/s5_pair_large_sample.json (office|D0=1, office|D0=10, office|D0=100: mean, se, R1bar_exact)
The mean-field $\Rz$ does not: at $\Dz=1$ the index case of the office infects $2.26\pm0.06$ people where
$\Rz=5.11$, and the ratio $\Rhat/\Rz$ is $0.44$, $0.45$, $0.45$ and $0.60$ in the office, supermarket, classroom
and metro carriage. It rises to $0.84$--$0.93$ at $\Dz=100$.
% src: data/bsc_sim/03_scenes.json (R_index_alone_mean, R0_ngm); s6_scenes_numbers.json: scenes.*.Rhat_over_R0

With the same-cell kernel the mean field is not adequate for the outbreak probability where cells are smaller
than the contact range: in the metro carriage
$P_{\mathrm{br}}-P_{\mathrm{maj}}$ is $0.54$ at $\Dz=10$ ($0.888$ against $0.349$; attack rate of major
outbreaks $17.3\pct$ against $100.0\pct$) and still $0.031$ at $\Dz=100$ ($98.6\pct$ against $100.0\pct$); in the
classroom it is $0.08$ at $\Dz=10$ ($89.1\pct$ against $98.8\pct$) and $0.014$ at $\Dz=100$
(Section~\ref{s6_scenes:sec:kernel}).
% src: data/bsc_sim/03_scenes.json: "metro|D0=10|samecell" (p_major 0.349, p_major_branching 0.888, attack_major_mean 0.173, ode_attack 1.000),
%      "metro|D0=100|samecell" (0.858, 0.889, 0.986), "classroom|D0=10|samecell" (0.695, 0.777, 0.891, 0.988), "classroom|D0=100|samecell" (0.771, 0.785)

Peak height and peak day are a different matter. At $\Dz=10$ the simulated peaks ($39.0\,\%$, $30.6\,\%$, $35.0\,\%$,
$27.5\,\%$) are lower and about four to five days later than those of the mean-field equations even when these
are started from a single cell ($44.1\,\%$, $33.1\,\%$, $41.3\,\%$, $36.3\,\%$). At $\Dz=100$ the two agree to within $2.5$
percentage points and about one day in the office, the supermarket and the classroom. In the metro carriage they
do not: its mobility is $0.028\,\Dz$ on area average (Table~\ref{M-s2_model:tab:scenes}), the infection has to
travel along the 45 cells of the carriage, and the mean-field peak depends on the seed ($64.6\,\%$ on day $7.0$
for a uniform seed, $42.2\,\%$ on day $13.8$ for a single cell) and stays above the simulated $35.3\,\%$ on day
$16.5$.
% src: data/checks/simulation.json: metro_area_mean_D_over_D0 (0.0282, recomputed from scenes/metro.json; notes/VERIFICATION.md: Section 4.2, simulations, item 14)
% src: data/bsc_sim/03_scenes.json (peak_prev_major_mean, peak_time_major_mean, ode_peak_prev, ode_peak_time);
%      data/s6_scenes_ode_pointseed.json (point_seed.peak_mean, day_mean; rows D0 = 10 and 100)

\paragraph{Replication.}
The second simulator gave $P_{\mathrm{maj}}=0.498$, $0.429$, $0.603$ and $0.864$ for the four scenes at
$\Dz=1$ (3{,}000--4{,}000 runs each) and $0.734$ for the office at $\Dz=10$, against $0.496$, $0.436$, $0.614$,
$0.879$ and $0.762$ here. The two largest differences are $1.6$ and $2.1$ standard errors. The rows for
$\Dz=100$, the metro carriage at $\Dz=10$ and the same-cell rows at $\Dz\ge10$ were not re-run with the second simulator. Sampling prevalence every $0.5$ day lowers the reported peaks by $0.3$
to $1.0$ percentage point relative to the exact maximum.
% src: code/bsc_sim/verify/v05_scenes.json (p_major, n; peak_exact_major - peak_grid_major);
%      s6_scenes_numbers.json: check_scenes.*.z_first_minus_check (1.56, 2.11), peak_exact_minus_grid_points (0.28..0.96);
%      notes/VERIFICATION.md: Section 4.2, simulations, summary (D0 = 100 rows not re-simulated)

\subsection{Outbreak probability and early extinction}\label{s6_scenes:sec:pmajor}

\begin{table}[tbp]
\centering
\caption{Offspring distribution of a lone index case and the probability of a major outbreak (1\,m kernel;
$\Dz$ in \cdunit{}). Columns 3--6: mean, variance, variance-to-mean ratio and probability of zero of the
number of secondary cases when only the uniformly placed index case transmits (40{,}000 runs per row). Columns
7--9: probability that a branching process does not die out, for a Poisson offspring law with the simulated
mean, for the mean-field process \eqref{M-s6_scenes:eq:branching}, and for a Galton--Watson process with the
simulated offspring law. Last column: fraction of 2{,}000 full epidemics with attack rate at least $10\,\%$
(Table~\ref{M-s6_scenes:tab:scenes}).}
\label{s6_scenes:tab:pmajor}
% src: data/bsc_sim/03b_outbreak_probability.json (experiments/03b_outbreak_probability.py)
\footnotesize
\setlength{\tabcolsep}{4pt}
\begin{tabular}{@{}lrcccccccc@{}}
\toprule
 & & \multicolumn{4}{c}{secondary cases of a lone index} & \multicolumn{3}{c}{branching processes} & simulation\\
\cmidrule(lr){3-6}\cmidrule(lr){7-9}\cmidrule(l){10-10}
Scene & $\Dz$ & mean & variance & ratio & $P(0)$ & Poisson & mean field & Galton--Watson & $P_{\mathrm{maj}}$ (95\,\% CI)\\
\midrule
% BEGIN AUTO s6_scenes:tab:pmajor
Office & 1 & 2.27 & 6.0 & 2.66 & 0.269 & 0.857 & 0.781 & 0.613 & 0.496 (0.474--0.518)\\
Supermarket & 1 & 2.00 & 5.3 & 2.68 & 0.310 & 0.796 & 0.667 & 0.531 & 0.436 (0.414--0.458)\\
Classroom & 1 & 2.58 & 6.7 & 2.60 & 0.228 & 0.903 & 0.767 & 0.688 & 0.614 (0.592--0.635)\\
Metro & 1 & 5.54 & 21.0 & 3.80 & 0.109 & 0.996 & 0.888 & 0.877 & 0.879 (0.864--0.893)\\
\addlinespace
Office & 10 & 3.82 & 16.3 & 4.25 & 0.194 & 0.976 & 0.783 & 0.757 & 0.762 (0.743--0.781)\\
Supermarket & 10 & 2.89 & 11.0 & 3.80 & 0.253 & 0.932 & 0.685 & 0.658 & 0.648 (0.627--0.669)\\
Classroom & 10 & 3.73 & 15.5 & 4.15 & 0.197 & 0.973 & 0.777 & 0.753 & 0.750 (0.731--0.768)\\
Metro & 10 & 6.35 & 31.9 & 5.02 & 0.104 & 0.998 & 0.888 & 0.883 & 0.878 (0.863--0.892)\\
% END AUTO s6_scenes:tab:pmajor
\bottomrule
\end{tabular}
\end{table}

Early extinction of a chain started by one case occurs in every scene: the fraction of runs below the
$10\,\%$ threshold is $0.50$, $0.56$, $0.39$ and $0.12$ at $\Dz=1$ and $0.24$, $0.35$, $0.25$ and $0.12$ at
$\Dz=10$ (office, supermarket, classroom, metro carriage).
% src: data/bsc_sim/03_scenes.json: 1 - p_major (s6_scenes_numbers.json: scenes.*.die_out)
Table~\ref{s6_scenes:tab:pmajor} shows what determines it. The number of secondary cases of an index case is
over-dispersed, with a variance $2.6$ to $5.0$ times the mean, and $10\,\%$ to $31\,\%$ of index cases infect no one.
% src: 03b_outbreak_probability.json: var_over_mean (2.598..5.022), p_zero (0.104..0.310)
An exponential infectious period alone produces this much: a Poisson number of infections over an exponentially
distributed period is geometric, with a variance-to-mean ratio of one plus the mean ($3.0$ to $7.4$ here). The
simulated ratios are lower, consistent with the limited number of people within reach.
% src: s6_scenes_numbers.json: pmajor_ranges.geometric_var_over_mean_min/max (2.9975, 7.3503)
Over-dispersion lowers the outbreak probability \citep{lloydsmith2005superspreading}: a Poisson law with the same
mean would give $0.80$--$1.00$. A Galton--Watson process fed with the simulated offspring law reproduces
$P_{\mathrm{maj}}$ to within $0.011$ wherever the attack-rate distribution is bimodal (all scenes at $\Dz=10$, the
metro carriage at $\Dz=1$).
% src: s6_scenes_numbers.json: pmajor_ranges.gw_minus_sim_other_rows_max_abs = 0.0105; 03b: p_major_poisson_same_mean (0.796..0.998)

In the other three scenes at $\Dz=1$ the comparison depends on the convention, because the notion of a major
outbreak is not sharp there. With a threshold of $5\,\%$, $10\,\%$ or $20\,\%$ the office gives $0.60$, $0.50$ or
$0.40$, the supermarket $0.49$, $0.44$ or $0.37$ and the classroom $0.70$, $0.61$ or $0.49$, whereas at $\Dz=10$
the three thresholds change no value by more than $0.014$.
% src: s6_scenes_numbers.json: threshold_sensitivity_D0_10_max_change = 0.0135
% src: s6_scenes_numbers.json: scenes.*.P_attack_ge_5pct, P_attack_ge_10pct, P_attack_ge_20pct
%      (computed from data/bsc_sim/03_curves_*_range1m.npz: attack)
The Galton--Watson values ($0.61$, $0.53$, $0.69$) are $0.07$--$0.12$ above the $10\,\%$ fractions and within
$0.05$ of the $5\,\%$ fractions, so no conclusion about later generations can be drawn from them. The mean-field
branching values ($0.78$, $0.67$, $0.77$) are too high under either convention, by $0.15$--$0.29$ and
$0.07$--$0.19$.
% src: s6_scenes_numbers.json: pmajor.*.gw_minus_sim (0.117, 0.095, 0.074), pmajor.*.branching_minus_sim (0.285, 0.231, 0.153);
%      differences to the 5 % fractions: pmajor.*.gw_minus_sim_5pct (0.018, 0.045, -0.011), pmajor.*.branching_minus_sim_5pct (0.186, 0.180, 0.067)

\subsection{Sensitivity to the contact kernel}\label{s6_scenes:sec:kernel}

The contact kernel is a modelling choice with large consequences when cells are smaller than the contact range.
With the strict same-cell kernel the metro carriage ($0.5$\,m cells) produces no outbreak at
$\Dz=1$: none of 2{,}000 runs reaches an attack rate of $10\,\%$ (upper $95\,\%$ limit $0.002$; none of 3{,}000 in the
second simulator) and the mean attack rate is $1.1\,\%$, although $\Rz=9.84$.
% src: 03_scenes.json: "metro|D0=1|samecell" (n_major = 0, p_major_hi = 0.0019, attack_all_mean = 0.0112, R0_ngm = 9.836);
%      code/bsc_sim/verify/v05_scenes.json: "metro|1.0|same" (n = 3000, p_major = 0)
The index case does infect: $\Ronebar=1.51$, $\Rhat=1.56\pm0.04$, and the pair-level matrix has
$\srad(\ngmone)=2.24$. But passengers hardly move at this mobility, the people the index case can reach are those
who share its cell, and the cases it makes are born in that cell with no one left to infect: over the 2{,}000
full epidemics the second generation is $0.61$ times as numerous as the first.
% src: 03_scenes.json: "metro|D0=1|samecell" (R1_pair_uniform_index = 1.514, R_index_alone_mean = 1.562, _se = 0.038,
%      R1_pair_ngm = 2.244, gen2_over_gen1 = 0.608); re-check: v05_scenes.json "metro|1.0|same".gen2_over_gen1 = 0.609
With the 1\,m kernel the same carriage is the most explosive scene ($P_{\mathrm{maj}}=0.879$). In the classroom
the same-cell rule gives $P_{\mathrm{maj}}=0.276$ against $0.614$ with the 1\,m kernel, with
$\srad(\ngmone)=2.27$ and a generation ratio of $0.90$. The difference between the kernels fades with mobility
and is gone in the classroom at $\Dz=100$ ($0.771$ against $0.779$), not yet in the metro carriage ($0.858$
against $0.891$).
% src: 03_scenes.json: "classroom|D0=1|samecell" (p_major = 0.2755, R1_pair_ngm = 2.265, gen2_over_gen1 = 0.899),
%      "classroom|D0=100|samecell|range1m" (0.7715, 0.779), "metro|D0=100|samecell|range1m" (0.858, 0.891)
These cases are the evidence for Remark~\ref{M-s5_pair:rem:threshold}: $\Ronebar>1$ and $\srad(\ngmone)>1$ do not
imply that outbreaks occur, and neither does $\Rz>1$. Nothing computed in this article predicts the threshold
of the discrete system at low mobility; it has to be simulated.

\subsection{Where infections happen}

\begin{table}[tbp]
\centering
\caption{Where infections happen: prediction against simulation (1\,m kernel; $\Dz$ in \cdunit{}).
(a)~Pearson correlation $r$, over the walkable cells, between the predicted and the simulated distribution of
the cell in which the cases of generation $g=1,2,3$ are infected (40{,}000 outbreaks per row, uniformly placed
index case, generations 0--2 infectious); mean-field prediction $\propto\ngm^{g}\one$, pair-level prediction
$\propto\ngmone^{g}\one$. `Index alone': first generation in 100{,}000 runs in which only the index case
transmits, with the value of $r$ expected from sampling noise if the pair-level prediction is exact. Last column:
cases of generation~2 per outbreak, simulated / pair level / mean field. (b)~Office scene: share of the
infections of generation $g$ that occur in a class of cells divided by the share of the walkable floor it
occupies (1 = uniform).}
\label{s6_scenes:tab:hotspot}
% src (a): data/bsc_sim/04_spatial.json (r_meanfield, r_pair, mean_cases_per_run, pred_cases_per_run_*),
%          data/bsc_sim/04b_first_generation.json (r_meanfield, r_pair, r_expected_if_exact)
% src (b): data/s6_scenes_numbers.json: office_cell_classes (code/s6_scenes_tables.py, from
%          data/bsc_sim/04_spatial_office_D1.npz and 04_spatial_office_D10.npz)
\footnotesize
\setlength{\tabcolsep}{3pt}
\begin{tabular}{@{}lrcccccc@{}}
\multicolumn{8}{@{}l}{(a) Correlation between predicted and simulated infection maps}\\
\toprule
 & & \multicolumn{2}{c}{$r$, generations 1 / 2 / 3} & \multicolumn{3}{c}{$r$, index alone} & cases per outbreak,\\
\cmidrule(lr){3-4}\cmidrule(lr){5-7}
Scene & $\Dz$ & mean field & pair level & mean field & pair level & if exact & generation 2\\
\midrule
% BEGIN AUTO s6_scenes:tab:hotspot:a
Office & 1 & 0.09 / $-$0.23 / $-$0.37 & 0.88 / 0.82 / 0.71 & 0.284 & 0.948 & 0.958 & 2.8 / 5.2 / 21.8\\
Supermarket & 1 & 0.37 / 0.15 / 0.22 & 0.91 / 0.85 / 0.85 & 0.550 & 0.981 & 0.977 & 2.6 / 4.0 / 11.6\\
Classroom & 1 & 0.76 / 0.67 / 0.68 & 0.96 / 0.93 / 0.91 & 0.844 & 0.989 & 0.989 & 2.8 / 6.8 / 23.6\\
Metro & 1 & 0.76 / 0.69 / 0.57 & 0.92 / 0.95 / 0.91 & 0.876 & 0.980 & 0.979 & 8.1 / 31.6 / 82.8\\
\addlinespace
Office & 10 & 0.77 / 0.52 / 0.41 & 0.92 / 0.85 / 0.81 & 0.915 & 0.985 & 0.985 & 8.1 / 14.4 / 21.6\\
Supermarket & 10 & 0.97 / 0.97 / 0.98 & 0.99 / 0.99 / 0.99 & 0.981 & 0.996 & 0.996 & 6.3 / 8.3 / 10.9\\
Classroom & 10 & 0.96 / 0.96 / 0.96 & 0.98 / 0.98 / 0.97 & 0.984 & 0.996 & 0.996 & 6.8 / 14.1 / 22.7\\
Metro & 10 & 0.77 / 0.67 / 0.60 & 0.93 / 0.95 / 0.93 & 0.889 & 0.982 & 0.984 & 10.3 / 40.3 / 82.6\\
% END AUTO s6_scenes:tab:hotspot:a
\bottomrule
\end{tabular}

\vspace{1.5ex}
\begin{tabular}{@{}lrrrccc@{}}
\multicolumn{7}{@{}l}{(b) Office scene: relative infection rate by class of cells, generations 1 / 2 / 3}\\
\toprule
Class of cells & $\Dz$ & cells & floor (\%) & simulation & pair level & mean field\\
\midrule
% BEGIN AUTO s6_scenes:tab:hotspot:b
Workstations bordering a corridor & 1 & 56 & 23.9 & 1.13 / 1.24 / 1.35 & 1.13 / 1.16 / 1.18 & 1.08 / 1.02 / 0.98\\
Other workstations & 1 & 100 & 42.7 & 0.94 / 0.89 / 0.88 & 0.94 / 0.91 / 0.89 & 1.08 / 1.11 / 1.12\\
Corridor & 1 & 39 & 16.7 & 0.97 / 1.09 / 1.16 & 0.92 / 0.93 / 0.93 & 0.62 / 0.55 / 0.52\\
Meeting room & 1 & 26 & 11.1 & 1.00 / 0.88 / 0.68 & 1.06 / 1.12 / 1.16 & 1.15 / 1.29 / 1.44\\
Kitchen & 1 & 13 & 5.6 & 0.94 / 0.76 / 0.58 & 1.00 / 0.98 / 0.96 & 0.92 / 0.83 / 0.75\\
\addlinespace
Workstations bordering a corridor & 10 & 56 & 23.9 & 1.15 / 1.19 / 1.20 & 1.13 / 1.12 / 1.12 & 1.08 / 1.05 / 1.04\\
Other workstations & 10 & 100 & 42.7 & 1.06 / 1.08 / 1.09 & 1.06 / 1.07 / 1.07 & 1.08 / 1.08 / 1.08\\
Corridor & 10 & 39 & 16.7 & 0.71 / 0.72 / 0.72 & 0.68 / 0.67 / 0.67 & 0.62 / 0.60 / 0.59\\
Meeting room & 10 & 26 & 11.1 & 0.93 / 0.81 / 0.76 & 1.04 / 1.05 / 1.05 & 1.15 / 1.24 / 1.30\\
Kitchen & 10 & 13 & 5.6 & 0.87 / 0.79 / 0.78 & 0.89 / 0.84 / 0.81 & 0.92 / 0.87 / 0.84\\
% END AUTO s6_scenes:tab:hotspot:b
\bottomrule
\end{tabular}
\end{table}

\begin{figure}[tbp]
\centering
\includegraphics[width=\textwidth]{fig5_2_spatial_pattern_en}
\caption{Where infections happen in the office scene at $\Dz=1\cellday$ ($\Npop=100$, 234 walkable cells, 1\,m
kernel, equal to the same-cell kernel here; 40{,}000 outbreaks from a uniformly placed index case, generations 0--2
infectious). (a)--(c)~Frequency with which the cases of generation~2 are infected in each cell, relative to a
uniform distribution: mean-field prediction $\propto\ngm^2\one$, pair-level prediction $\propto\ngmone^2\one$,
and simulation. Grey cells are barriers; thin lines separate the zones of Fig.~\ref{M-s2_model:fig:scenes} (the
cross is the corridor, the block at the top right the meeting room, the block at the bottom right the kitchen).
(d)~The same three maps cell by cell, with the correlation coefficients $r$ of
Table~\ref{s6_scenes:tab:hotspot}. (e)~Share of infections divided by share of floor for five classes of cells
and generations 1, 2, 3 (numbers in Table~\ref{s6_scenes:tab:hotspot}b).}
% src: code/s6_scenes_fig_spatial.py <- data/bsc_sim/04_spatial_office_D1.npz, 04_spatial.json,
%      data/s6_scenes_numbers.json (office_cell_classes); output figures/fig5_2_spatial_pattern_{en,zh}.pdf
\label{s6_scenes:fig:spatial}
\end{figure}

For a uniformly placed index case the cell in which a case of generation $g$ is infected has a distribution
proportional to $\ngm^{g}\one$ in the mean-field model and, in the pair-level description, to $\ngmone^{g}\one$.
Table~\ref{s6_scenes:tab:hotspot}(a) compares both with the infection trees of 40{,}000 outbreaks per scene.

In the office at $\Dz=1$ the mean-field map is uncorrelated with the simulated one in the first generation and
anti-correlated in the next two ($r=0.09$, $-0.23$, $-0.37$), whereas the pair-level map gives $r=0.88$, $0.82$,
$0.71$. The second simulator gives $0.11$, $-0.23$, $-0.37$ and $0.88$, $0.82$, $0.73$.
% src: data/bsc_sim/04_spatial.json: "office|D0=1".gen1..3 (r_meanfield 0.092, -0.235, -0.371; r_pair 0.875, 0.816, 0.709);
%      code/bsc_sim/verify/v07_hotspot.json: gen1..3 (r_mf 0.105, -0.234, -0.366; r_pair 0.878, 0.824, 0.726)
The pair-level map is exact for the first generation of a lone index case, and the simulation is compatible with
that: with only the index case transmitting, $r=0.948$ where $0.958$ is expected from sampling noise, and the
goodness-of-fit statistic is between $0.85$ and $1.20$ per degree of freedom in the eight rows (smallest
$p$-value $0.02$, office at $\Dz=1$).
% src: data/bsc_sim/04b_first_generation.json (r_pair, r_expected_if_exact, chi2_pair, dof);
%      s6_scenes_numbers.json: first_generation_chi2 (chi2_over_dof_min = 0.846, _max = 1.203; p = 0.0186)
For later generations it is an approximation, and it over-predicts their size, because a new case starts inside
the cluster of its infector: in generation~2 there are $2.8$ cases per outbreak in the simulation, $5.2$ in the
pair-level description and $21.8$ in the mean field.
% src: 04_spatial.json: "office|D0=1".gen2 (mean_cases_per_run = 2.769, pred_cases_per_run_pair = 5.167, pred_cases_per_run_mf = 21.76)
Over all scenes at $\Dz=1$ the pair-level correlation lies between $0.71$ and $0.96$ and the mean-field one between $-0.37$ and $0.76$.
At $\Dz=10$ the mean-field map is adequate in the supermarket and the classroom ($r=0.96$--$0.98$) but not in
the office ($0.77$, $0.52$, $0.41$) or the metro carriage ($0.77$, $0.67$, $0.60$).
% src: s6_scenes_numbers.json: hotspot_pair_r_range_D0_1 (0.709, 0.957), hotspot_meanfield_r_range_D0_1 (-0.371, 0.764);
%      04_spatial.json rows D0 = 10

Which cells are hot is shown in Fig.~\ref{s6_scenes:fig:spatial} and Table~\ref{s6_scenes:tab:hotspot}(b). The
mean-field map of the first generation is proportional to $q$: it marks the meeting room ($q=1.5$) and the
workstations ($q=1.4$) and leaves the corridor ($q=0.8$) cold, at $0.62$ of the uniform rate. In the simulation
at $\Dz=1$ the corridor is close to average. With only the index case transmitting it receives $15.4\,\%$ of the
infections on $16.7\,\%$ of the floor (pair level $15.3\,\%$, mean field $10.3\,\%$), and in full outbreaks its
relative rate is $0.97$, $1.09$ and $1.16$ in generations 1--3.
% src: code/bsc_sim/verify/v07_hotspot.json: zone_C (obs 0.1537, pair 0.1532, mf 0.1026, area 0.1667; 100 000 runs);
%      s6_scenes_numbers.json: office_cell_classes."D0=1".classes.C.sim (0.974, 1.087, 1.160), .meanfield[0] = 0.615
So the corridors are not the hotspots. The hottest cells are workstation cells that border a corridor: their
relative rate is $1.13$, $1.24$ and $1.35$, against $0.94$, $0.89$ and $0.88$ for the other workstation cells
(second simulator: $1.15$, $1.26$, $1.36$ and $0.94$, $0.88$, $0.87$), and in each generation at least 14 of
the 15 hottest cells are workstation cells, 13 to 15 of them next to a corridor.
% src: s6_scenes_numbers.json: office_cell_classes."D0=1".classes.W_border.sim, .W_interior.sim, .top15 (zones, n_border_workstations);
%      data/checks/simulation.json: office_hotspots_D0_1.relative_rate_generations_1_3 (values of the re-check;
%      notes/VERIFICATION.md: Section 4.2, simulations, item 6)
A plausible reason is that these cells combine a high $q$ with a faster exchange of neighbours: the hop rate
\eqref{M-s2_model:eq:hop} across a workstation--corridor edge is $0.5\,\Dz$ against $0.3\,\Dz$ between two
workstation cells. The meeting room, the hotspot of the mean-field analysis of Section~\ref{M-s4_geometry:sec},
falls from the average rate to $0.68$ by the third generation, where the mean field predicts $1.44$. This is
consistent with local depletion: only about 11 of the 100 people are in it at any time and, at this mobility,
few others arrive.
% src: office_cell_classes."D0=1".classes.M.sim (0.998, 0.877, 0.677), .meanfield (1.154, 1.295, 1.436);
%      hop rates: 2*0.3*1.5/(0.3+1.5) = 0.5 and 0.3 (code/bsc_sim/scenes/office.json: D_W = 0.3, D_C = 1.5); 26/234 x 100 = 11.1
At $\Dz=10$ the corridor is cold in the simulation as well ($0.71$--$0.72$; mean field $0.59$--$0.62$) and the
workstation cells are above average whether or not they border a corridor ($1.15$--$1.20$ and $1.06$--$1.09$).
% src: office_cell_classes."D0=10".classes (C.sim, C.meanfield, W_border.sim, W_interior.sim)
A statement such as `infection concentrates in the meeting room and at the workstations' is thus a mean-field
statement. For discrete people at low mobility the map has to be taken from $\ngmone$ or from simulation, and
all of this concerns a room whose occupants perform a random walk without home positions.

\subsection{Ranking of the scenes}\label{s6_scenes:sec:ranking}

At $\Dz=1$ the four scenes rank in the same order, metro carriage, classroom, office, supermarket, by $\Rz$
($9.24$, $5.87$, $5.11$, $4.39$), by $\Rhat$ ($5.59$, $2.62$, $2.26$, $1.97$), by $P_{\mathrm{maj}}$ and by the
attack rate over all runs.
% src: 03_scenes.json (R0_ngm, R_index_alone_mean, p_major, attack_all_mean at D0 = 1); s6_scenes_numbers.json: ranking_D0_1
The order is not robust. The attack rate of major outbreaks at $\Dz=1$ is lowest in the classroom ($39.2\,\%$),
and at $\Dz=10$ the office and the classroom cannot be told apart ($\Rhat=3.90\pm0.09$ and $3.72\pm0.09$;
$P_{\mathrm{maj}}=0.762$ and $0.750$).
% src: 03_scenes.json rows D0 = 10; s6_scenes_numbers.json: ranking_D0_1.by_attack_major, ranking_D0_10
Above all, the ranking belongs to the closed-room idealisation, in which the same people share the room for the
whole infectious period of $1/\gamma=7.1$ days. Nobody spends a week in a metro carriage. Counted per visit at
$\Dz=1$, an index case causes $0.175$ infections in an office day of 8 hours, $0.027$ in a one-hour lecture,
$0.018$ in a 20-minute metro ride and $0.008$ in a supermarket visit of 27 minutes
(Section~\ref{M-s6_scenes:sec:interventions}, Table~\ref{M-s7_interventions:tab:dwell}), which puts the office first and the
metro carriage third.
% src: data/bsc_sim/11_dwell_time.json: per_visit_sim (office 0.1753, classroom 0.0270, metro 0.0178, supermarket 0.0083),
%      visit_hours (8, 1, 0.333, 0.45)

\subsection{Heterogeneity of the infection efficiency}\label{s7_interventions:sec:hetero}

\begin{figure}[tbp]
  \centering
  \includegraphics[width=\textwidth]{s7_interventions_hetero_en}
  \caption{Two-zone room ($20\times13$ cells; workspace 70\,\% of the floor with mobility $0.3\Dz$, corridor
  30\,\% with $2\Dz$; same-cell kernel; FD) as a function of the contrast $c$ of
  \eqref{M-s7_interventions:eq:contrast}; the dotted vertical line marks $c=0.583$ ($q_{\mathrm W}=1.5$,
  $q_{\mathrm C}=0.5$). (a)~Mean-field $\Rz$ for four mobility scales; the grey band is the interval
  $[\beta\langle q\rangle/\gamma,\beta\qmax/\gamma]$ of Theorem~\ref{M-s3_r0:thm:bounds}.
  (b,\,c)~Counted secondary cases $\Rhat$ of the index case (only the index transmits; 95\,\% intervals;
  20,000 runs per point for $\Npop=100$, 4,000 for $\Npop=1000$) at $\Dz=1$ and $\Dz=10$. Filled symbols: index
  placed with the Perron vector of the pair-level matrix $\ngmone$, to be compared with $\srad(\ngmone)$ (solid
  lines). Open symbols: index placed uniformly, to be compared with $\Ronebar$ (dashed lines). `Uniform
  mobility' is the control with $D=0.81\Dz$ in every cell (Perron-placed index only).
  (d)~Mean attack rate over all runs with 95\pct{} intervals (2,000 epidemics per point), $\Dz=1$
  (filled, solid) and $\Dz=10$ (open, dashed). Numbers in Table~\ref{s7_interventions:tab:hetero}.}
  % src: code/s7_interventions_fig_hetero.py <- data/bsc_sim/05_heterogeneity.json, 05_heterogeneity_extra.json, data/s7_interventions_hetero_se.json (panel d)
  \label{s7_interventions:fig:hetero}
\end{figure}

The two-zone room of Fig.~\ref{s7_interventions:fig:hetero} has $20\times13$ cells without barriers, a lattice
spacing of 1.5\,m (so that the 1\,m kernel is the same-cell kernel), FD calibration, and $\Npop=100$ (0.38 other
people per cell) or $\Npop=1000$ (3.8). The contrast $c$ of \eqref{M-s7_interventions:eq:contrast} is $0$ in the
uniform room, $0.583$ for $q_{\mathrm W}=1.5$, $q_{\mathrm C}=0.5$, and negative when the higher efficiency is in
the corridor. Table~\ref{s7_interventions:tab:hetero} gives the numbers. The mean-field $\Rz$ is $4.972$ at
$c=0.583$ and $\Dz=1$, and $4.427$ ($+3.3\,\%$) at $\Dz=10$.
% src: code/bsc_sim/experiments/05_heterogeneity.py (set-up); data/bsc_sim/05_heterogeneity.json keys D0=1|N=100|c=0.5833, D0=10|N=100|c=0.5833 (R0_ngm)

\begin{table}[tbp]
  \centering\footnotesize\setlength{\tabcolsep}{3.5pt}
  \caption{Two-zone room of Fig.~\ref{s7_interventions:fig:hetero} (same-cell kernel, FD): mean-field $\Rz$,
  pair-level $\srad(\ngmone)$, counted secondary cases $\Rhat$ (95\pct{} interval; 20,000 runs for $\Npop=100$,
  4,000 for $\Npop=1000$) for an index placed with the Perron vector of $\ngmone$ and for a uniformly placed
  index, probability of a major outbreak and mean attack rate over all runs ($\pm$ standard error; 2,000
  epidemics per row). Mobility `zones': $0.3\Dz$ in the workspace and $2\Dz$ in the corridor; `uniform':
  $0.81\Dz$ everywhere; $c=-2$ is the hot-corridor control. $^{\dagger}$From the
  verification (20,000 runs; standard error 0.02 to 0.03); not simulated in the main runs.}
  \label{s7_interventions:tab:hetero}
  % src: data/s7_interventions_tables.tex (generated by code/s7_interventions_numbers.py) <- data/bsc_sim/05_heterogeneity.json, 05_heterogeneity_extra.json; dagger entries: code/bsc_sim/verify/v09_het.json (sim_uniform); last two columns: data/s7_interventions_hetero_se.json (stored|... for N = 100: the stored 2,000-run batches re-run from their seeds; new|... for N = 1000: 2,000 further epidemics, seed 5503, in place of the 300 stored ones)
  \begin{tabular}{@{}rrlccccccc@{}}
    \toprule
    $\Dz$ & $\Npop$ & mobility & $c$ & $\Rz$ & $\srad(\ngmone)$ & $\Rhat$, Perron index & $\Rhat$, uniform index & $P(\text{major})$ & attack \\
    \midrule
    1 & 100 & zones & 0.00 & 4.29 & 2.75 & 2.75 (2.71--2.79) & 2.42 (2.39--2.46) & $0.541\pm0.011$ & $0.347\pm0.008$ \\
     &  &  & 0.58 & 4.97 & 2.12 & 2.10 (2.07--2.13) & 2.07 (2.04--2.11) & $0.460\pm0.011$ & $0.231\pm0.006$ \\
     &  &  & 1.00 & 5.64 & 2.13 & 2.12 (2.09--2.15) & 1.69 (1.66--1.72) & $0.345\pm0.011$ & $0.119\pm0.004$ \\
    1 & 1000 & zones & 0.00 & 4.29 & 3.97 & 3.87 (3.74--3.99) & 3.97 (3.84--4.11) & $0.766\pm0.009$ & $0.740\pm0.009$ \\
     &  &  & 0.58 & 4.97 & 4.33 & 4.23 (4.09--4.37) & 3.79 (3.65--3.93) & $0.727\pm0.010$ & $0.688\pm0.009$ \\
     &  &  & 1.00 & 5.64 & 4.82 & 4.76 (4.60--4.92) & 3.57 (3.42--3.71) & $0.639\pm0.011$ & $0.573\pm0.010$ \\
    1 & 100 & uniform & 0.00 & 4.29 & 2.78 & 2.76 (2.72--2.80) & 2.76$^{\dagger}$ & $0.646\pm0.011$ & $0.503\pm0.009$ \\
     &  &  & 0.58 & 4.76 & 2.82 & 2.82 (2.78--2.86) & 2.61$^{\dagger}$ & $0.620\pm0.011$ & $0.442\pm0.008$ \\
    1 & 100 & zones & $-2.00$ & 9.14 & 4.97 & 5.02 (4.96--5.09) & 2.50$^{\dagger}$ & $0.381\pm0.011$ & $0.244\pm0.007$ \\
    \midrule
    10 & 100 & zones & 0.00 & 4.29 & 3.69 & 3.65 (3.59--3.70) & 3.69 (3.64--3.75) & $0.748\pm0.010$ & $0.729\pm0.009$ \\
     &  &  & 0.58 & 4.43 & 3.60 & 3.59 (3.54--3.64) & 3.52 (3.46--3.57) & $0.744\pm0.010$ & $0.717\pm0.009$ \\
     &  &  & 1.00 & 4.66 & 3.63 & 3.64 (3.58--3.69) & 3.37 (3.32--3.43) & $0.704\pm0.010$ & $0.671\pm0.010$ \\
    10 & 1000 & zones & 0.00 & 4.29 & 4.22 & 4.34 (4.19--4.49) & 4.06 (3.92--4.20) & $0.753\pm0.010$ & $0.742\pm0.009$ \\
     &  &  & 0.58 & 4.43 & 4.33 & 4.35 (4.21--4.50) & 4.08 (3.94--4.23) & $0.745\pm0.010$ & $0.730\pm0.010$ \\
    10 & 100 & uniform & 0.58 & 4.36 & 3.86 & 3.89 (3.84--3.95) & -- & $0.737\pm0.010$ & $0.720\pm0.009$ \\
    10 & 100 & zones & $-2.00$ & 5.70 & 4.93 & 4.89 (4.82--4.96) & -- & $0.614\pm0.011$ & $0.585\pm0.010$ \\
    \bottomrule
  \end{tabular}
\end{table}

\paragraph{Discrete individuals, $\Dz=1$, $\Npop=100$.} Everything that is counted \emph{falls} when the contrast
is switched on (Fig.~\ref{s7_interventions:fig:hetero}b,d; numbers in Table~\ref{s7_interventions:tab:hetero}): $\Rhat$ of
an index case placed with the Perron vector of $\ngmone$ goes from $2.75$ to $2.10$ ($-24\,\%$), $\Rhat$ of a
uniformly placed index from $2.42$ to $2.07$ ($-15\,\%$), the ratio of second- to first-generation cases
from $2.19$ to $1.58$ ($-28\,\%$), the probability of a major outbreak from $0.541$ to $0.460$ and the mean
attack rate from $0.347$ to $0.231$ ($-33\,\%$).
% src: data/bsc_sim/05_heterogeneity.json keys D0=1|N=100|c=0.0000 and c=0.5833 (sim_R_perron, sim_R_uniform, sim_gen2_over_gen1, p_major, attack_all); data/s7_interventions_numbers.json hetero/changes_D1_N100
(The second simulator gives $2.77\pm0.02\to2.11\pm0.02$ and $0.341\to0.239$ for the first
and the last of these.)
% src: code/bsc_sim/verify/v09_het.json keys <zone variant>|D0=1.0|N=100|c=0.000, c=0.583
The cause is not the heterogeneity of $q$ as such but its coincidence with low mobility: this layout puts
the high efficiency where people hardly move, where an infective keeps meeting the same neighbours and its
contacts fall on people who are already infected (Section~\ref{M-s5_pair:sec}). Two controls confirm this. With
the same contrast and \emph{uniform} mobility ($0.81\,\Dz$ everywhere, the area mean) $\Rhat$ rises slightly,
$2.76\to2.82$; with the high efficiency in the \emph{mobile} corridor ($c=-2$) $\Rhat$ of a Perron-placed
index reaches $5.02$.
% src: data/bsc_sim/05_heterogeneity_extra.json keys extra|uniformD|D0=1|N=100|c=0.0000, c=0.5833; extra|<zone variant>|D0=1|N=100|c=-2.0000 (sim_R_perron)
In both controls the mean attack rate nevertheless falls ($0.503\to0.442$ and $0.347\to0.244$): a low-$q$
zone is a place where an index case tends to fizzle and which the epidemic reaches late.
% src: data/bsc_sim/05_heterogeneity_extra.json (attack_all), same keys

\paragraph{High occupancy.} With $\Npop=1000$ the pair-level $\srad(\ngmone)$ rises with the contrast
($3.97$, $4.33$, $4.82$ at $c=0$, $0.58$, $1$), and so does $\Rhat$ of a Perron-placed index ($3.87$, $4.23$,
$4.76$; re-check $3.92$, $4.33$, $4.78$, each $\pm0.08$ to $0.09$).
% src: data/bsc_sim/05_heterogeneity.json keys D0=1|N=1000|c=... (R1_pair_ngm, sim_R_perron); code/bsc_sim/verify/v09_het.json keys <zone variant>|D0=1.0|N=1000|c=... (sim_perron, se)
This is the only sense in which the simulation shows the rise of $\Rz$. For a uniformly placed index the
count falls in these runs ($3.97$, $3.79$, $3.57$; re-check $3.96$, $3.86$, $3.82$). It is bounded above by
$\beta\langle q\rangle/\gamma=4.29$ for every contrast (Proposition~\ref{M-s5_pair:prop:saturation}), so it
cannot follow the rise of $\Rz$; but it is not monotone in the contrast in general: with the high efficiency
in the mobile corridor ($c<0$, $\Npop=100$) $\Ronebar$ rises from $2.41$ at $c=0$ to $2.66$ at $c=-1$ before
falling to $2.44$ at $c=-2$ ($\Dz=1$), and from $3.69$ to $3.81$ at $\Dz=10$.
% src: data/bsc_sim/05_heterogeneity.json keys D0=1|N=1000|c=... (sim_R_uniform; 4000 counted runs); code/bsc_sim/verify/v09_het.json (sim_uniform)
% src: data/bsc_sim/05_heterogeneity.json D0=1|N=100|c=0.0000 (R1_pair_uniform 2.411), D0=10 (3.688); data/bsc_sim/05_heterogeneity_extra.json extra|<zone variant>|D0=1|N=100|c=-0.5 ... -2.0 (R1_pair_uniform 2.589, 2.662, 2.623, 2.441), D0=10|c=-1.0 (3.808)
The outbreak probability and the mean attack rate fall as well: in 2{,}000 epidemics per contrast they are
$0.766\pm0.009$, $0.727\pm0.010$, $0.639\pm0.011$ and $0.740\pm0.009$, $0.688\pm0.009$, $0.573\pm0.010$
(standard errors). The first sample of 300 epidemics per contrast ($0.737$, $0.690$, $0.630$ and $0.712$,
$0.654$, $0.564$, each $\pm0.025$ to $0.028$) resolved only the difference between $c=0$ and $c=1$. $\Rz$ is
the growth factor per generation once cases are distributed as the incidence map $\wvec$, not the fate of a
typical introduction.
% src: data/s7_interventions_hetero_se.json keys new|D0=1|N=1000|c=0.0000, 0.5833, 1.0000 (p_major, p_major_se, attack_all, attack_all_se; n = 2000, seed 5503); stored|D0=1|N=1000|c=... (n = 300; p_major_se 0.025-0.028, attack_all_se 0.025; reproduces_stored true); code/s7_interventions_hetero_se.py

\paragraph{$\Dz=10$.} For $\Npop=100$ the contrast $c=0.583$ changes $\Rhat$ of a Perron-placed index from $3.65$ to
$3.59$, the outbreak probability from $0.748$ to $0.744$ and the mean attack rate from $0.729$ to $0.717$
(all less than 2\pct, and within one standard error of $0.010$ for the last two); $\Rhat$ of a uniformly
placed index falls by 5\pct{} ($3.69\to3.52$) (Fig.~\ref{s7_interventions:fig:hetero}c,d). For $\Npop=1000$
the outbreak probability is $0.753\pm0.010$ and $0.745\pm0.010$ at $c=0$ and $0.58$, and the mean attack rate
$0.742\pm0.009$ and $0.730\pm0.010$: no change is resolved.
% src: data/bsc_sim/05_heterogeneity.json keys D0=10|N=100|c=0.0000, c=0.5833; data/s7_interventions_numbers.json hetero/changes_D10_N100; data/s7_interventions_hetero_se.json keys stored|D0=10|N=100|c=... (p_major_se, attack_all_se), new|D0=10|N=1000|c=0.0000, 0.5833

\subsection{Barriers with identical starts}\label{s7_interventions:sec:barriers}

To isolate the effect of barriers we use a uniform room ($20\times13$ cells, $\Npop=100$, $q=1.2$, so
$\beta q/\gamma=4.29$, uniform mobility $D$, same-cell kernel) and close randomly chosen cells at a nominal
density $\rho_B$. The barrier sets of one layout are nested; walkable pockets cut off by barriers are closed
as well, so that the realised density at nominal $0.30$ is $0.35$. In all arms of a layout the same people
start on the same cells (drawn from the cells that are walkable in the densest arm), the index case is the
same person on the centre cell, and each replicate uses the same random seed: the arms differ only by the
barriers. Fig.~\ref{s7_interventions:fig:barriers} shows a scan over seven
densities at $D=1$; Table~\ref{s7_interventions:tab:barriers} gives the
effects for 40 further layouts at $D=1$ and $D=10$, with standard errors across layouts, and for the
re-check's 40 layouts.
% src: code/bsc_sim/experiments/06_obstacles.py (design); code/s7_interventions_barriers.py; code/bsc_sim/verify/v08b_obstacles_many.py; realised density: data/s7_interventions_barriers.json D=1/arms/rhoB=0.30|FD/realised = 0.352

\begin{enumerate}[label=(\roman*),leftmargin=*,itemsep=2pt,topsep=3pt]
\item \emph{Mean field.} Under FD, $\Rz=\beta q/\gamma=4.29$ at every barrier density
  (Theorem~\ref{M-s4_geometry:thm:closed}). Under DD it rises as the floor shrinks,
  by \eqref{M-s7_interventions:eq:dd}: $4.76$ at 10\,\% and $6.7$ at nominal 30\,\%.
  % src: data/s7_interventions_barriers.json D=1/arms/rhoB=0.10|DD, rhoB=0.30|DD (R0); data/bsc_sim/06_obstacles.json (R0_ngm)
\item \emph{Pair level.} Under FD the pair-level $\Rone$ of the centre index falls a little, $3.02\to2.94\to2.80$
  at $\rho_B=0$, $0.10$, $0.30$ for $D=1$ and $3.90\to3.87\to3.75$ for $D=10$: barriers lengthen the time two
  people remain neighbours.
\item \emph{Outcomes at $D=1$, FD.} The mean attack rate falls by $0.040\pm0.003$ at 10\,\% barriers and by
  $0.223\pm0.011$ at 30\,\% (re-check: $0.047\pm0.004$ and $0.235\pm0.011$), the peak prevalence of major
  outbreaks from 27\pct{} to 25\pct{} and 18\pct, while the outbreak probability falls much
  less ($0.70\to0.69\to0.62$). Barriers act by stopping outbreaks part-way on a fragmented floor, not by
  preventing take-off.
\item \emph{Peak time at $D=1$, FD.} The peak of major outbreaks is delayed by $0.81\pm0.12$ day at 10\,\%
  barriers (re-check: $0.82\pm0.10$) and by $0.93\pm0.23$ day at 30\,\% (re-check: $0.71\pm0.30$): a real delay
  of about 0.8 day. The scan of
  Fig.~\ref{s7_interventions:fig:barriers}c cannot resolve it (its across-layout intervals are $\pm1.1$ day or
  wider).
\item \emph{$D=10$, FD.} Barriers at 10\,\% leave the attack rate unchanged ($+0.005\pm0.005$) and at 30\,\%
  lower it by $0.025\pm0.005$; the peak is lower by 0.8 and 5.3 percentage points and later by $0.34\pm0.05$
  and $2.05\pm0.16$ days.
\item \emph{DD.} Barriers also crowd people together. At $D=1$ the two effects cancel at 10\,\%
  ($-0.001\pm0.004$) and fragmentation wins at 30\,\% ($-0.085\pm0.015$); at $D=10$ crowding wins and the
  attack rate \emph{rises}, by $0.026\pm0.005$ and $0.077\pm0.005$, with the peak 1.0 and 2.8 days earlier.
\end{enumerate}
% src (items ii-vi): data/s7_interventions_barriers.json keys D=1, D=10: arms/*/R1_centre, p_major, peak_major; paired/*/attack, peak_major, peak_day_major (mean, SE across 40 layouts); re-check: code/bsc_sim/verify/v08b_obstacles_many.json (diff_0.1_ar, diff_0.3_ar, diff_0.1_pkt, diff_0.3_pkt); +-1.1 day: data/bsc_sim/06_obstacles.json (peak_time_major_mean, 95 % half-width over 8 layouts), data/s7_interventions_numbers.json barriers/first_8x500
Two caveats. Starting everyone on cells that are walkable in the densest arm is not the stationary law of
the sparser arms; the re-check measured mean attack rates of $0.600$ and $0.589$ at zero barriers for the two
ways of starting, a difference within noise.
% src: notes/VERIFICATION.md, Section 4.2, simulations, further point 7
And although the densities stay below the blocked-site percolation threshold of the square lattice, $0.407$
\citep{stauffer1994percolation}, about 5\,\% of this small floor is already cut off as pockets at nominal
$0.30$.

\begin{figure}[tbp]
  \centering
  \includegraphics[width=0.86\textwidth]{fig6_3_obstacle_scan_en}
  \caption{Barrier scan with identical starts: uniform room ($20\times13$ cells, $\Npop=100$, $q=1.2$,
  $D=1\cellday$, same-cell kernel), random nested barriers at nominal density $\rho_B$; 8 layouts with 500
  epidemics per layout and density; bars are 95\,\% intervals across layouts. Blue: FD; orange: DD (pair hazard
  frozen at its barrier-free value).
  (a)~Mean-field $\Rz$ (dashed), pair-level $\Rone(x_c)$ of the centre index (solid)
  and counted $\Rhat$ of the index (open circles, 500 runs per layout with only the index transmitting).
  (b)~Peak prevalence and (c)~peak day of major outbreaks. (d)~Mean attack rate over all runs. Effects with
  layout-level standard errors, including the peak delay that panel~(c) cannot resolve, are in
  Table~\ref{s7_interventions:tab:barriers}.}
  % src: figures/fig6_3_obstacle_scan_en.pdf (code/bsc_sim/experiments/21_figures_more.py <- data/bsc_sim/06_obstacles.json)
  \label{s7_interventions:fig:barriers}
\end{figure}

\begin{table}[tbp]
  \centering\footnotesize\setlength{\tabcolsep}{3.5pt}
  \caption{Barriers with identical starts in the uniform room of Fig.~\ref{s7_interventions:fig:barriers}:
  means over 40 nested layouts (1,000 epidemics per layout and arm at $D=1$, 500 at $D=10$). $\rho_B$: nominal
  barrier density; $\Rone(x_c)$: pair-level number of the centre index; $P$: probability of a major outbreak;
  attack: mean attack rate over all runs; peak and day: peak prevalence (\%) and peak day of major outbreaks.
  The $\Delta$ columns are differences from the barrier-free arm of the same layout, mean $\pm$ standard error
  across layouts ($\Delta$peak in percentage points, $\Delta$day in days). $^{\ddagger}$Re-check:
  40 other layouts, 1,000 epidemics each, second simulator.}
  \label{s7_interventions:tab:barriers}
  % src: data/s7_interventions_tables.tex (code/s7_interventions_numbers.py) <- data/s7_interventions_barriers.json (code/s7_interventions_barriers.py); ddagger rows: code/bsc_sim/verify/v08b_obstacles_many.json
  \begin{tabular}{@{}cccccccccccc@{}}
    \toprule
    $D$ & law & $\rho_B$ & $\Rz$ & $\Rone(x_c)$ & $P$ & attack & peak & day & $\Delta$attack & $\Delta$peak & $\Delta$day \\
    \midrule
    1 & FD & 0.00 & 4.29 & 3.02 & 0.698 & 0.609 & 27.2 & 21.5 & -- & -- & -- \\
    1 & FD & 0.10 & 4.29 & 2.94 & 0.688 & 0.569 & 24.7 & 22.3 & $-0.040\pm0.003$ & $-2.5\pm0.1$ & $+0.81\pm0.12$ \\
    1 & FD & 0.30 & 4.29 & 2.80 & 0.625 & 0.386 & 17.6 & 22.4 & $-0.223\pm0.011$ & $-9.6\pm0.4$ & $+0.93\pm0.23$ \\
    1 & DD & 0.10 & 4.76 & 3.16 & 0.713 & 0.608 & 26.4 & 21.2 & $-0.001\pm0.004$ & $-0.8\pm0.1$ & $-0.33\pm0.12$ \\
    1 & DD & 0.30 & 6.68 & 3.69 & 0.727 & 0.524 & 22.0 & 20.3 & $-0.085\pm0.015$ & $-5.2\pm0.5$ & $-1.22\pm0.25$ \\
    1 & FD$^{\ddagger}$ & 0.10 & 4.29 & -- & 0.683 & 0.564 & 24.6 & 22.3 & $-0.047\pm0.004$ & $-2.4\pm0.1$ & $+0.82\pm0.10$ \\
    1 & FD$^{\ddagger}$ & 0.30 & 4.29 & -- & 0.627 & 0.376 & 17.0 & 22.2 & $-0.235\pm0.011$ & $-10.0\pm0.3$ & $+0.71\pm0.30$ \\
    \midrule
    10 & FD & 0.00 & 4.29 & 3.90 & 0.757 & 0.744 & 44.0 & 14.3 & -- & -- & -- \\
    10 & FD & 0.10 & 4.29 & 3.87 & 0.764 & 0.749 & 43.2 & 14.7 & $+0.005\pm0.005$ & $-0.8\pm0.1$ & $+0.34\pm0.05$ \\
    10 & FD & 0.30 & 4.29 & 3.75 & 0.743 & 0.719 & 38.7 & 16.4 & $-0.025\pm0.005$ & $-5.3\pm0.3$ & $+2.05\pm0.16$ \\
    10 & DD & 0.10 & 4.76 & 4.25 & 0.779 & 0.770 & 46.2 & 13.3 & $+0.026\pm0.005$ & $+2.2\pm0.1$ & $-1.05\pm0.05$ \\
    10 & DD & 0.30 & 6.68 & 5.47 & 0.826 & 0.821 & 48.8 & 11.5 & $+0.077\pm0.005$ & $+4.8\pm0.6$ & $-2.78\pm0.21$ \\
    \bottomrule
  \end{tabular}
\end{table}

\subsection{Occupancy}\label{s7_interventions:sec:occupancy}

In the office scene we vary the number of people from $\Npop=20$ to $400$ (0.03 to 0.68 persons per m$^2$;
2,000 epidemics per point; Fig.~\ref{s7_interventions:fig:occupancy}).
% src: data/bsc_sim/07_occupancy.json (density_per_m2, n_rep)
 Under FD the mean-field $\Rz$ does not
depend on $\Npop$ at all: it is $5.11$ ($\Dz=1$) and $4.68$ ($\Dz=10$) for every occupancy. Discrete people do
feel occupancy, through saturation: with fewer people an infective exhausts its neighbours sooner. At $\Dz=1$
the counted $\Rhat$ is $1.47\pm0.04$ for $\Npop=50$, $2.29\pm0.05$ for $100$ and $3.70\pm0.09$ for $400$, following
$\Ronebar=1.52$, $2.27$, $3.65$; the attack rate of major outbreaks is 27\,\%, 46\,\% and 90\,\%. At $\Dz=10$
the dependence is weaker ($\Rhat=3.2$, $3.8$, $4.5$). The elasticity of $\Ronebar$ with respect to $\Npop$
near $\Npop=100$ (between 75 and 150) is $+0.48$ at $\Dz=1$ and $+0.17$ at $\Dz=10$, against $0$ for $\Rz$.
% src: data/bsc_sim/07_occupancy.json keys D0=1|N=50|FD, N=100, N=400, D0=10|... (R0_ngm, R1_pair_uniform, R_index_alone, R_index_alone_se, attack_major_mean); elasticities: data/s7_interventions_numbers.json occupancy/elasticity|D0=1|FD, D0=10|FD (pair)
Under DD (reference $\Npop=100$), $\Rz$ is proportional to $\Npop-1$ by \eqref{M-s7_interventions:eq:dd}, from
$0.98$ at $\Npop=20$ to $20.6$ at $\Npop=400$ for $\Dz=1$, and the elasticity is $+1.0$ at both mobilities.
Halving the office from 100 to 50 people then takes $\Ronebar$ from $2.27$ to $1.12$ and the outbreak
probability from $0.49$ to $0.27$ at $\Dz=1$ ($3.79\to1.88$ and $0.74\to0.53$ at $\Dz=10$), whereas under FD it
takes them to $1.52$ and $0.40$ ($3.22$ and $0.72$).
% src: data/bsc_sim/07_occupancy.json keys D0=1|N=20|DD, N=400|DD, N=50|DD, N=50|FD, N=100|FD and the D0=10 analogues (R0_ngm, R1_pair_uniform, p_major)

Whether a capacity limit works is therefore an \emph{assumption about contact behaviour} --- do people keep
the same number of contacts in an emptier room? --- and not an output of the model.

\begin{figure}[tbp]
  \centering
  \includegraphics[width=\textwidth]{fig6_4_occupancy_en}
  \caption{Occupancy dependence in the office scene (1\,m kernel, 2,000 epidemics per point), $\Dz=1$ (top) and
  $\Dz=10$ (bottom). Blue: FD; orange: DD with the pair hazard frozen at its value for $\Npop=100$.
  (a,\,d)~Mean-field $\Rz$ (dashed), pair-level $\Ronebar$ (solid) and counted
  $\Rhat$ (open circles, 95\pct{} intervals). (b,\,e)~Probability of a major outbreak: simulation (dots, 95\,\%
  intervals) and mean-field branching value (dashed); for $\Npop<50$ a `major' outbreak means two or three
  cases. (c,\,f)~Peak prevalence of major outbreaks: simulation (dots) and lattice mean-field equations
  \eqref{M-s2_model:eq:meanfield} (dashed).}
  % src: figures/fig6_4_occupancy_en.pdf (code/bsc_sim/experiments/21_figures_more.py <- data/bsc_sim/07_occupancy.json)
  \label{s7_interventions:fig:occupancy}
\end{figure}

\subsection{Interventions}\label{s7_interventions:sec:measures}

Table~\ref{s7_interventions:tab:interventions} applies five measures to the office scene. Two of them are
\emph{inputs}, not results: `ventilation' halves $q$ in every cell and `masks' multiply $\beta$ by $0.3$, by
assumption.

\emph{Uniform measures} scale $\Rz$ exactly ($5.11\to2.55$ and $1.53$ at $\Dz=1$), but the pair-level number
less than proportionally ($\Ronebar=2.27\to1.51$ and $1.05$), because a lower hazard also wastes fewer
contacts; the outbreak probability falls from $0.49$ to $0.30$ and $0.15$. \emph{Halving the occupancy}
leaves $\Rz$ unchanged under FD and halves it under DD (Section~\ref{s7_interventions:sec:occupancy}); in
either case the peak comes earlier, on day 13 instead of day 24 at $\Dz=1$, because a smaller group is
exhausted sooner. \emph{Partitions} (39 workstation cells turned
into barriers) leave $\Rz$ at $5.11$ under FD and raise it to $6.14$ under DD at $\Dz=1$; the mean attack rate
falls from $0.232$ to $0.166$ (FD) or $0.195$ (DD). At $\Dz=10$ it goes from $0.737$ to $0.693$ (FD) or $0.751$
(DD); the decrease of $\Rz$ under FD there ($4.68\to4.63$, $-1.1\pct$) arises because the cells removed have
an above-average $q$ (Section~\ref{M-s4_geometry:sec:examples}).
% src: data/bsc_sim/08_interventions.json keys office|D0=1|... and office|D0=10|... (R0_ngm, R1_pair_uniform, p_major, peak_time_major_mean, attack_all_mean)

\emph{Combination} (ventilation, half occupancy and partitions together). At $\Dz=1$, $\Rz=2.56$ (FD) or $1.52$
(DD), $\Ronebar=1.10$ or $0.82$ and $\Rhat=1.09\pm0.03$ or $0.78\pm0.02$; at $\Dz=10$, $\Ronebar=1.79$ or
$1.16$. In the classroom the combination gives $\Ronebar=1.21$ (FD) or $0.86$ (DD) at $\Dz=1$ and $1.70$ or
$1.10$ at $\Dz=10$.
% src: data/bsc_sim/08_interventions.json keys office|D0=1|combination|FD, |DD, office|D0=10|combination|FD, |DD, classroom|D0=1|combination|FD, |DD, classroom|D0=10|combination|FD, |DD
A pair-level number below one is reached only for discrete people at low mobility under the DD law, and it
must not be read as a control criterion: in that arm $\Rz$ is $1.52$ and 16.5\,\% of the runs still reach an
attack rate of 10\,\% (five of 50 people), and conversely a pair-level number above one does not imply
outbreaks (Remark~\ref{M-s5_pair:rem:threshold}).
% src: data/bsc_sim/08_interventions.json key office|D0=1|combination|DD (p_major = 0.165)

\emph{Targeted ventilation.} Halving $q$ on the quarter of the cells with the largest elasticity $e_x$
(Theorem~\ref{M-s4_geometry:thm:hotspots}) lowers $\Rz$ more than treating the quarter with the smallest ($4.67$
against $5.11$ at $\Dz=1$; $4.10$ against $4.37$ at $\Dz=10$). For discrete people at $\Dz=1$ it makes no
measurable difference which quarter is treated: $\Ronebar=2.08$, $2.09$ and $2.06$ and mean attack rates
$0.195$, $0.205$ and $0.188$ (standard errors about $0.005$) for the hottest quarter, the coldest quarter and the hottest
quarter by the pair-level elasticity; at $\Dz=10$ the attack rates are $0.675$, $0.696$ and $0.682$.
% src: data/bsc_sim/08_interventions.json keys office|D0=1|hotspot ventilation (25 % of cells)|FD, coldspot ..., pair-level hotspot ..., and D0=10 (R0_ngm, R1_pair_uniform, attack_all_mean, attack_all_sd/sqrt(2000))

\begin{table}[tbp]
  \centering\footnotesize\setlength{\tabcolsep}{4pt}
  \caption{Interventions in the office scene (1\,m kernel; 2,000 epidemics per arm; reference occupancy for DD:
  $\Npop=100$ in the unmodified scene). $\Rhat$: counted secondary cases of the index ($\pm$ standard error);
  $P$: probability of a major outbreak; attack: mean attack rate over all runs; peak and day: peak prevalence
  (\%) and peak day of major outbreaks. Ventilation and masks are assumptions on $q$ and $\beta$, not model
  outputs. Partitions: 39 workstation cells turned into barriers. Combination: ventilation, half occupancy and
  partitions. A value $\Ronebar<1$ is not a control criterion (see text).}
  \label{s7_interventions:tab:interventions}
  % src: data/s7_interventions_tables.tex (code/s7_interventions_numbers.py) <- data/bsc_sim/08_interventions.json
  \begin{tabular}{@{}cllcccccccc@{}}
    \toprule
    $\Dz$ & measure & law & $\Npop$ & $\Rz$ & $\Ronebar$ & $\Rhat$ & $P$ & attack & peak & day \\
    \midrule
    1 & baseline & FD & 100 & 5.11 & 2.27 & $2.30\pm0.05$ & 0.491 & 0.232 & 12.8 & 24.3 \\
    1 & ventilation, $q\to q/2$ & FD & 100 & 2.55 & 1.51 & $1.55\pm0.04$ & 0.295 & 0.102 & 9.0 & 23.3 \\
    1 & masks, $\beta\to0.3\beta$ & FD & 100 & 1.53 & 1.05 & $1.05\pm0.03$ & 0.147 & 0.048 & 7.1 & 21.8 \\
    1 & capacity 50\,\% & FD & 50 & 5.11 & 1.52 & $1.51\pm0.04$ & 0.393 & 0.126 & 11.7 & 13.0 \\
    1 & capacity 50\,\% & DD & 50 & 2.53 & 1.12 & $1.13\pm0.03$ & 0.258 & 0.080 & 10.2 & 13.9 \\
    1 & partitions & FD & 100 & 5.11 & 2.21 & $2.20\pm0.05$ & 0.459 & 0.166 & 10.7 & 20.3 \\
    1 & partitions & DD & 100 & 6.14 & 2.42 & $2.36\pm0.06$ & 0.506 & 0.195 & 11.4 & 19.9 \\
    1 & combination & FD & 50 & 2.56 & 1.10 & $1.09\pm0.03$ & 0.274 & 0.077 & 9.9 & 12.2 \\
    1 & combination & DD & 50 & 1.52 & 0.82 & $0.78\pm0.02$ & 0.165 & 0.055 & 9.0 & 12.3 \\
    \midrule
    10 & baseline & FD & 100 & 4.68 & 3.79 & $3.76\pm0.09$ & 0.760 & 0.737 & 39.0 & 16.7 \\
    10 & ventilation, $q\to q/2$ & FD & 100 & 2.34 & 2.09 & $2.09\pm0.06$ & 0.522 & 0.404 & 21.2 & 27.6 \\
    10 & masks, $\beta\to0.3\beta$ & FD & 100 & 1.40 & 1.30 & $1.30\pm0.04$ & 0.290 & 0.131 & 10.9 & 27.8 \\
    10 & capacity 50\,\% & FD & 50 & 4.68 & 3.22 & $3.20\pm0.07$ & 0.697 & 0.629 & 36.2 & 15.8 \\
    10 & capacity 50\,\% & DD & 50 & 2.31 & 1.88 & $1.85\pm0.05$ & 0.533 & 0.345 & 21.6 & 19.3 \\
    10 & partitions & FD & 100 & 4.63 & 3.60 & $3.60\pm0.09$ & 0.730 & 0.693 & 34.2 & 19.6 \\
    10 & partitions & DD & 100 & 5.55 & 4.14 & $4.13\pm0.10$ & 0.777 & 0.751 & 38.0 & 17.0 \\
    10 & combination & FD & 50 & 2.31 & 1.79 & $1.74\pm0.05$ & 0.504 & 0.285 & 18.8 & 19.5 \\
    10 & combination & DD & 50 & 1.37 & 1.16 & $1.19\pm0.04$ & 0.329 & 0.123 & 12.6 & 16.7 \\
    \bottomrule
  \end{tabular}
\end{table}

\subsection{Schedules and dwell time}\label{s7_interventions:sec:time}

\begin{proposition}[Duty cycle]\label{s7_interventions:prop:duty}
  Let $0<f\le1$ and let Assumption~\ref{M-s2_model:ass:A} hold. Suppose that movement and transmission operate
  only during the occupied intervals $[k,k+f)$, $k=0,1,2,\dots$ (time in days), and that recovery proceeds at all
  times: in \eqref{M-s2_model:eq:meanfield}, $\gen$ and $\Fm$ are multiplied by $\chi(t)$, equal to $1$ on the
  occupied intervals and to $0$ otherwise. Put $\theta(t)=\int_0^t\chi(s)\,\dd s$ and
  \begin{equation}\label{s7_interventions:eq:dutydetail}
    \Jm_f=\gen+\Fm-\frac{\gamma}{f}\,\Id,\qquad \Rz^{(f)}=\srad\Bigl(\Fm\bigl(\tfrac{\gamma}{f}\,\Id-\gen\bigr)^{-1}\Bigr).
  \end{equation}
  \begin{enumerate}
    \item[(a)] The linearised model $\dot I=[\chi(t)(\gen+\Fm)-\gamma\Id]\,I$ has the solution
      $I(t)=\exp[\theta(t)(\gen+\Fm)-\gamma t\,\Id]\,I(0)$; in particular $I(k)=\e^{kf\Jm_f}I(0)$ for integer $k$.
      Solutions of the nonlinear model satisfy $I(t)\le\exp[\theta(t)(\gen+\Fm)-\gamma t\,\Id]\,I(0)$ componentwise.
    \item[(b)] The growth factor per day is $\e^{f\sabs(\Jm_f)}$, and $\sabs(\Jm_f)$ has the sign of $\Rz^{(f)}-1$.
    \item[(c)] If $\gen=\Dz\genz$ and $\Rz(\Dz)$ denotes the value for continuous occupancy, then
      $\Rz^{(f)}(\Dz)=f\,\Rz(f\Dz)$. For the same-cell kernel,
      $f\,\Rz(\Dz)\le\Rz^{(f)}(\Dz)\le f\beta\qmax/\gamma$.
  \end{enumerate}
\end{proposition}
\begin{proof}
  (a) The matrices $A(t)=\chi(t)(\gen+\Fm)-\gamma\Id$ are constant on each occupied and each empty interval and
  commute with one another for all times, so the product of the exponentials over successive intervals is the
  exponential of $\int_0^tA(s)\,\dd s=\theta(t)(\gen+\Fm)-\gamma t\,\Id$; at $t=k$, $\theta=kf$. In the nonlinear
  model $S_y\le N^*_y$ gives $\dot I\le A(t)I$ componentwise, and since $A(t)$ is Metzler the comparison
  argument of Theorem~\ref{M-s3_r0:thm:ngm}(d) applies on each interval.
  (b) $\gen+\Fm$ is Metzler and irreducible, so $\e^{f\Jm_f}$ is a positive matrix with spectral radius
  $\e^{f\sabs(\Jm_f)}$ (Lemma~\ref{s3_r0:lem:facts}, P1); the sign statement is Theorem~\ref{M-s3_r0:thm:ngm}(c) with
  $\gamma$ replaced by $\gamma/f$.
  (c) $\Fm\bigl(\tfrac{\gamma}{f}\Id-\Dz\genz\bigr)^{-1}=f\,\Fm(\gamma\Id-f\Dz\genz)^{-1}$. The inequalities follow
  from Theorem~\ref{M-s3_r0:thm:mobility}(c), because $f\Dz\le\Dz$, and from Theorem~\ref{M-s3_r0:thm:bounds}.
\end{proof}

\begin{remark}[What is exact and what is approximate]\label{s7_interventions:rem:duty}
  $\Rz^{(f)}$ is an exact threshold parameter of the periodic mean-field model, for any cycle length, because
  one and the same operator is switched on and off; in general periodic models the threshold is given by a
  next-generation operator on periodic functions and not by time-averaged parameters
  \citep{bacaer2006epidemic,wang2008threshold}. By part~(c) a duty cycle does more than scale $\Rz$ by $f$: it
  also lowers the effective mobility to $f\Dz$.
  As a \emph{number of secondary cases}, however, the substitution $\gamma\to\gamma/f$ is an approximation. It
  replaces the expected infectious time spent in the room by $f/\gamma$ ($2.381$ days for $f=1/3$), whereas a
  case that becomes infectious at the start of an occupied interval spends
  $(1-\e^{-\gamma f})/[\gamma(1-\e^{-\gamma})]=2.493$ days there ($+4.7\,\%$), and one that becomes infectious
  at a uniformly distributed time within an occupied interval, as secondary cases approximately do (their
  infection times within an interval follow the incidence, which changes during the interval), $2.383$ days
  ($+0.07\,\%$).
  % src: data/s7_interventions_schedule.json keys duty|office|D0=1 (in_room_time_rule, in_room_time_start, in_room_time_uniform); closed forms in code/s7_interventions_schedule.py
  For discrete individuals the same substitution in the pair-level formulas of Section~\ref{M-s5_pair:sec} is
  likewise an approximation. The exact pair-level value follows from the argument of
  Proposition~\ref{M-s5_pair:prop:pair}. Let $A=\gamma\Id+\mathbf{H}-\mathcal{M}^{\T}$ be the pair operator of
  that proposition and $u_0(z)$ the probability that the given person escapes infection when the index case
  becomes infectious at the start of an occupied interval with the pair in state $z$. During the occupied
  interval of length $f$ the pair moves and the hazard acts, so $(\e^{-At}\one)(z)$ is the probability that
  neither recovery nor infection has occurred by time $t\le f$; during the rest of the day nobody moves, the
  index recovers with probability $c_f=1-\e^{-\gamma(1-f)}$, and otherwise the next day starts from the same
  state. Hence
  \begin{equation}\label{s7_interventions:eq:dutyexact}
    u_0=\gamma A^{-1}\bigl(\one-\e^{-Af}\one\bigr)+\e^{-Af}\bigl(c_f\,\one+(1-c_f)\,u_0\bigr),
  \end{equation}
  a linear fixed-point equation whose map is a contraction in the maximum norm (factor at most
  $\e^{-\gamma}<1$, because $\e^{-Af}$ is non-negative with row sums at most $\e^{-\gamma f}$), solved by
  iteration; the expected number of secondary cases is $(\Npop-1)$ times the mean of $1-u_0$ over the
  $\ncell^2$ pair states. For an index case that becomes infectious at a time uniformly distributed over the
  occupied interval, the first day is shortened to its remaining part $r$ and the same expression is averaged
  over $r\in(0,f)$. With $f=1$ the iteration reproduces the solution of \eqref{M-s5_pair:eq:pair} to
  $5\times10^{-10}$.
  % src: code/s7_interventions_schedule.py (duty_cycle_exact); data/s7_interventions_schedule.json key duty_exact|control_f1|office|D0=1 (abs_diff 4.6e-10), duty_exact|office|D0=1 (fixed_point_residual 9e-13, 152 iterations)
\end{remark}

\paragraph{Office, eight hours a day ($f=1/3$).} The mean-field number drops from $5.11$ to
$\Rz^{(f)}=1.75$ at $\Dz=1$ and from $4.68$ to $1.60$ at $\Dz=10$; the pair-level number with
$\gamma\to\gamma/f$ is $\Ronebar=0.889$ and $1.335$ (Table~\ref{M-s7_interventions:tab:dwell}a). The exact
pair-level values of \eqref{s7_interventions:eq:dutyexact} are $0.930$ at $\Dz=1$ and $1.397$ at $\Dz=10$ for
an index case that becomes infectious at the start of a working day, $4.6\pct$ and $4.7\pct$ above the rule,
as its longer time in the room leads one to expect, and $0.889$ and $1.336$ for one that becomes infectious
at a uniformly distributed time of the working day, $0.07\pct$ above the rule.
% src: data/s7_interventions_schedule.json keys duty_exact|office|D0=1, duty_exact|office|D0=10 (exact_start 0.9296, 1.3975; exact_uniform_phase 0.8892, 1.3359; rule_gamma_over_f 0.8885, 1.3349; start_over_rule_pct 4.62, 4.69; uniform_over_rule_pct 0.07), duty|office|D0=1, D0=10 (R0_duty)
Simulation with the schedule imposed explicitly agrees: $0.9309\pm0.0009$ and $0.8889\pm0.0009$ at $\Dz=1$
(2{,}000{,}000 runs each) and $1.400\pm0.002$ and $1.337\pm0.002$ at $\Dz=10$ (800{,}000 runs each), all
within $1.6$ standard errors of the exact values.
% src: data/s7_interventions_schedule.json keys duty_large|office|D0=1, duty_large|office|D0=10 (count_start, count_start_se, count_uni, count_uni_se, z_start 1.52 and 1.26, z_uni -0.33 and 0.82); smaller samples are listed below
In full epidemics with the schedule, 6.1\pct{} of the runs (95\pct{} interval 5.1--7.2\pct) reach a 10\pct{} attack
rate at $\Dz=1$ and 30.2\pct{} (28.2--32.2\pct) at $\Dz=10$. The index case of an office occupied eight hours a
day thus infects about $0.9$ people on average at $\Dz=1$ and $1.3$ to $1.4$ at $\Dz=10$, against $2.3$ and
$3.8$ for a room occupied around the clock; this is not a threshold statement
(Remark~\ref{M-s5_pair:rem:threshold}).
% src: data/bsc_sim/11_dwell_time.json keys office|D0=1, office|D0=10 (p_major, p_major_lo, p_major_hi)
For a class that meets one hour a day ($f=1/24$) the numbers are $\Rz^{(f)}=0.27$, $\Ronebar=0.18$ with the
rule and $0.194$ exactly (index infectious from the start of the lecture; counted $0.189\pm0.003$) at
$\Dz=1$, and no run in 2,000 reaches a 10\pct{} attack rate: with $\beta=0.5\perday$ one daily lecture cannot
sustain transmission in this model.
% src: data/bsc_sim/11_dwell_time.json key classroom|D0=1 (R0_duty_cycle, R1_pair_duty_cycle, sim_R_schedule, p_major); data/s7_interventions_schedule.json key duty_exact|classroom|D0=1 (exact_start 0.1944)

\paragraph{Further simulation estimates of the duty-cycle counts.} Before the exact values of
\eqref{s7_interventions:eq:dutyexact} were computed, the number of secondary cases of an index case in the
office occupied eight hours a day was estimated by simulation several times. For an index case infectious
from the start of a working day at $\Dz=1$: $0.918\pm0.009$ (one batch of 20,000 runs), $0.944\pm0.006$ (the
second simulator), $0.928\pm0.003$ (pooled, in the re-check's report) and $0.935\pm0.003$ (200,000
runs), against the exact $0.930$; at $\Dz=10$, $1.408\pm0.012$ and $1.395\pm0.004$ against $1.397$. For an
index case infectious from a uniformly distributed time: $0.896\pm0.003$ and $1.331\pm0.004$ (200,000 runs
each) against $0.889$ and $1.336$. The 200,000-run samples at $\Dz=1$ lie $2.0$ and $2.5$ standard errors
above the exact values; the samples of 2,000,000 runs quoted in Section~\ref{s7_interventions:sec:time} do
not ($+1.5$ and $-0.3$).
% src: data/bsc_sim/11_dwell_time.json keys office|D0=1, office|D0=10 (sim_R_schedule and its interval); code/bsc_sim/verify/v11.json key office_duty_D1; data/checks/simulation.json: duty_cycle_office_D0_1.count_explicit_schedule_pooled (0.928 +- 0.003), .first_analysis_batch (0.918) (notes/VERIFICATION.md Section 4.2, simulations item 11); data/s7_interventions_schedule.json keys duty|office|D0=1, D0=10 (count_start, count_uniform_phase and their se), duty_exact|office|D0=* (z_mc_start 1.97, z_mc_uniform_phase 2.54), duty_large|office|D0=1 (z_start 1.52, z_uni -0.33)

\paragraph{One visit.} Table~\ref{M-s7_interventions:tab:dwell}b gives the expected number of people that one
infective infects during a single stay of length $T$. For a uniformly placed infective the mean-field value is
$\beta\langle q\rangle(1-\e^{-\gamma T})/\gamma$ whatever the mobility and the kernel, because the walk is
stationary; the counted exposure agrees with it within 0.3\,\%, and the counted infections are lower only
where pairs saturate within the visit. The values are $0.21$ for an office day ($0.175$ counted at $\Dz=1$,
$0.198$ at $\Dz=10$), $0.027$ for a lecture, $0.018$ for a 20-minute metro ride and $0.0086$ for a 27-minute
supermarket visit.
% src: data/bsc_sim/11_dwell_time.json keys <scene>|D0=1, <scene>|D0=10 (per_visit_meanfield, per_visit_sim, per_visit_exposure_sim; 200,000 runs each)
A per-visit number is not the reproduction number of a setting that is visited repeatedly: an infective who
rides the metro twice a day for the whole infectious period infects $f\beta\langle q\rangle/\gamma=0.25$
fellow passengers in expectation, and one who shops once a day $0.06$ fellow customers. Weighting each
setting by the fraction of time spent in it is the residence-time (Lagrangian) formulation of multi-setting
epidemic models \citep{bichara2015sis}; here it is applied to one setting at a time. These numbers, and
not the closed-room values, are what the settings contribute, and they undercut the ranking of
Section~\ref{M-s6_scenes:sec}: as closed rooms the scenes rank metro, classroom, office, supermarket
($\Rz=9.2$, $5.9$, $5.1$, $4.4$ at $\Dz=1$); with the dwell times above the mean-field numbers are $1.75$ for
the office and $0.27$ for the lecture ($\Rz^{(f)}$, closed groups), $0.25$ for the metro carriage and $0.06$
for the supermarket ($f\beta\langle q\rangle/\gamma$).
% src: data/bsc_sim/11_dwell_time.json (wellmixed_duty_cycle: metro 0.2515, supermarket 0.0612; R0_continuous; R0_duty_cycle)

\paragraph{A daily schedule of the parameters.} In the metro scene we let the mobility field and a transmission
multiplier follow a daily signal with two rush hours (24 hourly segments; the multiplier is $1$ in the rush
hours and $0.8$ otherwise, mean $0.833$), and compare three arms: constant rush-hour values (the scene of
Section~\ref{M-s6_scenes:sec}), the hourly schedule, and constants equal to the time averages of the mobility
field and of the multiplier. The number of people is $\Npop=310$ in all arms: only mobility and the multiplier
vary, not the occupancy.
% src: code/bsc_sim/experiments/09_metro_time.py; data/bsc_sim/09_metro_time.json key schedule (m_mean)

\begin{observation}[Time average]\label{s7_interventions:obs:average}
  In the metro scene the hourly schedule and its time average give the same outbreak probability, attack
  rate, peak prevalence and peak day within Monte Carlo error at $\Dz=1$, $10$ and $100$
  (Table~\ref{s7_interventions:tab:schedule}); for example the peak prevalence of
  major outbreaks is 23.0\pct{} against 22.9\pct{} on day 25.3 against 25.2 at $\Dz=1$.
\end{observation}
% src: data/bsc_sim/09_metro_time.json keys D0=1|varying, D0=1|mean, D0=100|...; data/s7_interventions_schedule.json keys metro|D0=10|peak, varying, mean
A test in the re-check, run with the production simulator, in the office scene (half of each day with hop rates multiplied by $1.8$
and transmission by $1.3$, the other half by $0.2$ and $0.7$, against constants; 6,000 epidemics each) agrees:
outbreak probability $0.491$ against $0.498$, peak prevalence 12.8\,\% against 12.8\,\%, peak day $23.7$ against
$23.3$ (standard error $0.3$).
% src: code/bsc_sim/verify/v14_schedule.json keys check_varying, check_mean
This is not a theorem: the hourly generators do not commute and the hop rates are not small compared with one
cycle per day. Relative to constant rush-hour values the schedule lowers the multiplier but raises the
mobility; the first effect wins at $\Dz=1$ (peak 25.0\,\% $\to$ 23.0\,\%, three days later), the second at
$\Dz=100$ (35.4\,\% $\to$ 44.2\,\%, 2.6 days earlier), and at $\Dz=10$ they nearly cancel.

\begin{table}[tbp]
  \centering\footnotesize\setlength{\tabcolsep}{5pt}
  \caption{Daily schedule in the metro scene ($\Npop=310$ in every arm, 1\,m kernel, FD, 2,000 epidemics per
  arm). Arms: constant rush-hour values (`peak'), hourly schedule of mobility and transmission multiplier
  (`schedule'), constants equal to the time averages (`average'). $\Rz$ is not defined by a single matrix
  for the schedule arm. 95\,\% intervals in parentheses.}
  \label{s7_interventions:tab:schedule}
  % src: data/s7_interventions_tables.tex (code/s7_interventions_numbers.py) <- data/bsc_sim/09_metro_time.json (D0 = 1, 100); data/s7_interventions_schedule.json (D0 = 10)
  \begin{tabular}{@{}clccccc@{}}
    \toprule
    $\Dz$ & arm & $\Rz$ & $P(\text{major})$ & attack, major (\%) & peak prevalence, major (\%) & peak day, major \\
    \midrule
    1 & peak & 9.24 & 0.871 (0.855--0.885) & 93.6 & 25.0 (24.7--25.2) & 22.2 (21.6--22.7) \\
    1 & schedule & -- & 0.837 (0.820--0.853) & 93.4 & 23.0 (22.7--23.2) & 25.3 (24.6--26.0) \\
    1 & average & 7.66 & 0.844 (0.827--0.859) & 93.7 & 22.9 (22.7--23.1) & 25.2 (24.5--25.9) \\
    \midrule
    10 & peak & 9.16 & 0.887 (0.873--0.901) & 97.9 & 27.8 (27.5--28.1) & 21.0 (20.5--21.5) \\
    10 & schedule & -- & 0.869 (0.854--0.883) & 99.3 & 28.6 (28.3--28.9) & 21.0 (20.5--21.4) \\
    10 & average & 7.58 & 0.862 (0.846--0.876) & 99.3 & 28.6 (28.4--28.9) & 20.9 (20.5--21.4) \\
    \midrule
    100 & peak & 9.08 & 0.880 (0.866--0.894) & 99.9 & 35.4 (35.1--35.8) & 16.5 (16.2--16.8) \\
    100 & schedule & -- & 0.853 (0.836--0.867) & 99.8 & 44.2 (43.9--44.6) & 13.9 (13.7--14.1) \\
    100 & average & 7.55 & 0.855 (0.839--0.870) & 99.8 & 44.4 (44.0--44.7) & 13.9 (13.8--14.1) \\
    \bottomrule
  \end{tabular}
\end{table}

\subsection{What remains at higher mobility}\label{s7_interventions:sec:regime}

The experiments above were run at $\Dz=1$ and $10$. The pair-level number $\Ronebar$ is exact for the first
generation at any mobility (Definition~\ref{M-s5_pair:def:R1}) and costs one linear solve, so each effect can be
followed up to $\Dz=2.4\times10^{3}\cellday$, the estimate of Section~\ref{M-s2_model:sec:regime} for people who
walk 5\pct{} of the time (Table~\ref{s7_interventions:tab:regime}). The effects
of the contrast in $q$ and of
barriers under FD shrink below 1\,\% at $\Dz=100$ and below 0.05\,\% at $\Dz=2.4\times10^{3}$. Simulation at
$\Dz=100$ agrees: the contrast $c=0.583$ changes the mean attack rate from $0.744$ to $0.750$ (2,000 epidemics each,
standard error $0.009$), and barriers at 10\,\% and 30\,\% change it by $+0.007\pm0.005$ and $+0.013\pm0.007$
(20 layouts with 250 epidemics each, FD).
% src: data/s7_interventions_regime.json keys det/* (R0, R1bar) and sim/twozone|D0=100|c=0.0000, c=0.5833 (attack_all_mean, attack_all_sd/sqrt(2000)), sim/barriers|D0=100|FD/paired (attack); code/s7_interventions_regime.py
Three things survive. Uniform reductions of $q$ or $\beta$ act on $\Ronebar$ almost in proportion, as they do
on $\Rz$. Occupancy matters in full under DD; under FD a residue of about 5\,\% remains for halving the office,
which is a finite-population effect (the limit $D\to\infty$ in Proposition~\ref{M-s5_pair:prop:torus}) and not a
spatial one. And the
time spent in the room enters through Proposition~\ref{s7_interventions:prop:duty} at any mobility. Epidemics
at $\Dz=2.4\times10^{3}$ were not simulated.

\begin{table}[tbp]
  \centering\footnotesize\setlength{\tabcolsep}{5pt}
  \caption{Relative change (\%) of deterministic reproduction numbers caused by the measures of this section,
  as a function of the mobility scale $\Dz$ (\cdunit). Two-zone room and barrier room: $\Npop=100$,
  same-cell kernel (barriers: mean over the 40 layouts of Table~\ref{s7_interventions:tab:barriers}); office
  scene: 1\,m kernel; DD reference: the unmodified room. Last row: the ratio $\Ronebar/\Rz$ in the office scene.}
  \label{s7_interventions:tab:regime}
  % src: data/s7_interventions_tables.tex (code/s7_interventions_numbers.py) <- data/s7_interventions_regime.json key det
  \begin{tabular}{@{}llcccc@{}}
    \toprule
    measure & quantity & $\Dz=1$ & $10$ & $100$ & $2.4\times10^{3}$ \\
    \midrule
    two-zone contrast $c=0.58$ & $\Rz$ & $+16.0$ & $+3.3$ & $+0.4$ & $+0.01$ \\
     & $\Ronebar$ & $-14.6$ & $-4.3$ & $-0.6$ & $-0.03$ \\
    barriers, $\rho_B=0.30$, FD & $\Ronebar$ & $-7.6$ & $-4.7$ & $-0.9$ & $-0.04$ \\
    barriers, $\rho_B=0.30$, DD & $\Ronebar$ & $+20$ & $+38$ & $+50$ & $+52$ \\
    capacity 50\,\%, FD (office) & $\Ronebar$ & $-32.9$ & $-15.2$ & $-6.2$ & $-4.5$ \\
    capacity 50\,\%, DD (office) & $\Ronebar$ & $-50.5$ & $-50.5$ & $-50.5$ & $-50.5$ \\
    ventilation, $q\to q/2$ (office) & $\Ronebar$ & $-33.4$ & $-45.0$ & $-48.3$ & $-48.8$ \\
    \midrule
    office scene, baseline & $\Ronebar/\Rz$ & 0.44 & 0.81 & 0.93 & 0.95 \\
    \bottomrule
  \end{tabular}
\end{table}

\subsection{Computational cost}\label{appD_tables:sec:cost}\label{s7_interventions:sec:cost}

Table~\ref{s7_interventions:tab:cost} compares, for the
office scene, two ways of computing \emph{the same quantity}, charging Monte Carlo only for the runs needed to
reach a stated accuracy (its cost scales as accuracy$^{-2}$); times are single-core CPU times on one laptop.
% src: code/bsc_sim/experiments/10_benchmark.py (method)

One sparse factorisation of $\gamma\Id-\gen$ followed by power iteration gives $\Rz$ in 1.4 to 4.6\,ms (within
$10^{-8}$ of a dense eigen-solve). Estimating the same number to 1\,\% by a particle power method costs 8 to 9
times more at $\Dz=1$ and 113 to 133 times more at $\Dz=10$ with a vectorised NumPy estimator, and 3 and 86
times more with a compiled one; at 0.1\,\% all factors are 100 times larger. The equal-accuracy speed-up is
thus a few to about $10^2$ at 1\,\% accuracy and a few hundred to about $10^4$ at 0.1\,\%; it grows with
mobility (more hops per infectious period) and depends on the implementation by a factor of two to three.
% src: data/bsc_sim/10_benchmark.json keys D0=1, D0=10 / T1_R0_meanfield (cpu_sparse_s, abs_err_sparse, speedup_at_1pct, speedup_at_0p1pct); code/bsc_sim/verify/v12_benchmark_rerun.json (same keys; re-run of the same script); code/bsc_sim/verify/v13_bench.json (speedup_1pct, speedup_01pct; compiled estimator)
The pair-lattice solve for $\Ronebar$ (conjugate gradients on $\ncell^2=54{,}756$ unknowns, 0.07 to 0.14\,s) is
about 20 ($\Dz=1$) to 90 ($\Dz=10$) times faster than counting secondary cases to 1\,\% and agrees with the
count within 0.1 standard error.
% src: data/bsc_sim/10_benchmark.json T3_R1_discrete (cpu_pair_s, speedup_at_1pct, theory_minus_mc_in_se); code/bsc_sim/verify/v12_benchmark_rerun.json T3_R1_discrete
For whole-epidemic outputs the mean-field route is not a substitute for simulation at $\Dz=1$: the branching
value of the outbreak probability is too high by $0.28$ and the attack rate of the lattice equations by $0.55$;
no speed-up is meaningful for a biased answer. At $\Dz=10$ both biases are $0.03$, and the lattice equations
(1.5\,s) are slower than simulating to 1\,\% (0.19\,s, 33 runs). The exact simulator takes 0.5\,ms for one
office epidemic at $\Dz=1$ (5,900 events) and 5.9\,ms at $\Dz=10$.
% src: data/bsc_sim/10_benchmark.json T4_p_major (theory_bias), T5_attack_rate (theory_bias, cpu_ode_s, mc_cpu_for_1pct_s, mc_runs_for_1pct), c_core_cpu_per_run_s, events_per_run
In short, the spectral formula and the pair solve are the inexpensive ways to obtain $\Rz$ and $\Ronebar$,
and simulation is cheap enough to be the method of choice for outbreak probabilities and sizes.

\begin{table}[tbp]
  \centering\footnotesize\setlength{\tabcolsep}{4pt}
  \caption{Cost of deterministic evaluation against Monte Carlo for the same quantity (office scene,
  $\Npop=100$; single-core CPU). $\noff(x)$: offspring number \eqref{M-s3_r0:eq:offspring} of a case introduced in one
  cell; CG: conjugate gradients. Speed-up: Monte Carlo CPU time needed for a relative standard error of 1\,\%
  (an absolute standard error of 0.01 for the outbreak probability) divided by the deterministic CPU time; the
  three values for $\Rz$ are the first run, the re-check's re-run of the same script, and the
  re-check's compiled estimator; the two values for the other rows are the first run and the re-check's
  re-run; no speed-up is quoted where the deterministic value is biased. Bias: deterministic value minus
  simulated value of the individual-based model.}
  \label{s7_interventions:tab:cost}
  % src: data/bsc_sim/10_benchmark.json; code/bsc_sim/verify/v12_benchmark_rerun.json; code/bsc_sim/verify/v13_bench.json (collected in data/s7_interventions_numbers.json, key cost)
  \begin{tabular}{@{}llccccc@{}}
    \toprule
     & & \multicolumn{2}{c}{deterministic CPU} & \multicolumn{2}{c}{speed-up at equal accuracy} & bias \\
    \cmidrule(lr){3-4}\cmidrule(lr){5-6}
    quantity & deterministic method & $\Dz=1$ & $\Dz=10$ & $\Dz=1$ & $\Dz=10$ & $\Dz=1$; $10$ \\
    \midrule
    $\Rz=\srad(\ngm)$ & sparse LU, power iteration & 4.6\,ms & 1.4\,ms & 8; 9; 3 & 133; 113; 86 & none \\
    $\noff(x)$, one cell & one sparse solve & 0.27\,ms & 0.29\,ms & 35; 22 & 331; 351 & none \\
    $\Ronebar$ & CG on the pair lattice & 68\,ms & 137\,ms & 18; 22 & 85; 97 & none \\
    $P(\text{major})$ & mean-field branching & 9.4\,ms & 9.7\,ms & -- & -- & $+0.28$; $+0.03$ \\
    attack rate, major & lattice equations \eqref{M-s2_model:eq:meanfield} & 0.40\,s & 1.54\,s & -- & -- & $+0.55$; $+0.03$ \\
    \bottomrule
  \end{tabular}
\end{table}
